\documentclass[11pt,a4paper]{article}
\usepackage[italian]{babel}
\usepackage[margin=2.5cm]{geometry}
\usepackage{parskip}
\usepackage{amsmath, amssymb, amsthm}
\usepackage{booktabs}
\usepackage{array}
\usepackage{xcolor}
\usepackage{needspace}
\usepackage{enumitem}
\usepackage{listings}
\usepackage{tikz}
\usepackage[most]{tcolorbox}
\usepackage[hidelinks]{hyperref}
\usetikzlibrary{decorations.pathreplacing, shapes.geometric, arrows.meta, positioning}

% Niente vedove/orfane: i paragrafi non lasciano righe isolate a fine/inizio pagina
\widowpenalty=10000
\clubpenalty=10000

% Elenchi più compatti
\setlist{itemsep=2pt, parsep=0pt, topsep=4pt, partopsep=0pt}

% --- Riquadri con bordo nero, senza numero (si spezzano tra due pagine) ---
% Uso: \begin{definizione} ... \end{definizione}
%      \begin{teorema}[Nome] ... \end{teorema}  (il nome è facoltativo)
\tcbset{nota/.style={enhanced jigsaw, breakable, colback=white,
  colframe=black, boxrule=0.5pt, sharp corners}}
\NewTColorBox{definizione}{}{nota,
  before upper={\textbf{Definizione.}\ }}
\NewTColorBox{teorema}{o}{nota,
  before upper={\textbf{Teorema\IfValueT{#1}{ (#1)}.}\ }}
\NewTColorBox{proposizione}{o}{nota,
  before upper={\textbf{Proposizione\IfValueT{#1}{ (#1)}.}\ }}
\NewTColorBox{proprieta}{o}{nota,
  before upper={\textbf{Proprietà\IfValueT{#1}{ (#1)}.}\ }}
\NewTColorBox{assioma}{o}{nota, boxrule=1pt,
  before upper={\textbf{Assioma\IfValueT{#1}{ (#1)}.}\ }}

% --- Ambienti semplici (senza riquadro) ---
% Il titolo sta in cima al blocco, anche se il contenuto inizia con un riquadro o
% con codice; e non resta mai da solo in fondo a una pagina.
% Uso: \begin{esempio}[Titolo facoltativo] ... \end{esempio}
\newcommand{\newrunin}[2]{%
  \NewDocumentEnvironment{#1}{o}{%
    \par\Needspace{4\baselineskip}\addvspace{\medskipamount}\noindent
    \textbf{#2}\IfValueT{##1}{ (##1)}\textbf{.}\ \ignorespaces}%
  {\par\addvspace{\medskipamount}}}
\newrunin{esempio}{Esempio}
\newrunin{esercizio}{Esercizio}
\newrunin{osservazione}{Osservazione}

% V e F nelle tavole di verità
\newcommand{\V}{\textsf{V}}
\newcommand{\F}{\textsf{F}}

% --- Codice C ---
% Uso: \begin{codice} ... \end{codice}      codice C numerato
%      \begin{comando} ... \end{comando}    comandi da terminale
%      \begin{risultato} ... \end{risultato} output di un programma
%      \lstinline|printf("ciao");|           codice dentro al testo
\lstdefinestyle{c}{
  language=C,
  basicstyle=\ttfamily\small,
  keywordstyle=\color{blue}\bfseries,
  commentstyle=\color{gray}\itshape,
  stringstyle=\color{red!70!black},
  numbers=left,
  numberstyle=\tiny\color{gray},
  numbersep=8pt,
  columns=flexible,
  keepspaces=true,
  breaklines=true,
  showstringspaces=false,
  tabsize=4,
  morekeywords={include, define},
  % lettere accentate nei commenti
  literate={à}{{\`a}}1 {è}{{\`e}}1 {é}{{\'e}}1 {ì}{{\`i}}1
           {ò}{{\`o}}1 {ù}{{\`u}}1 {È}{{\`E}}1 {À}{{\`A}}1
}
\lstset{style=c}
\newtcblisting{codice}{listing only, nota, left=6mm,
  listing options={style=c}}
\newtcblisting{comando}{listing only, nota,
  listing options={basicstyle=\ttfamily\small, numbers=none, language={},
    columns=flexible, keepspaces=true, breaklines=true}}
\newtcblisting{risultato}{listing only, enhanced jigsaw, breakable,
  colback=black!5, colframe=black, boxrule=0.5pt, sharp corners,
  listing options={basicstyle=\ttfamily\small, numbers=none, language={},
    columns=flexible, keepspaces=true, breaklines=true}}

% --- Diagrammi di flusso (simboli standard) ---
\tikzset{
  terminale/.style = {ellipse, draw, minimum width=2.5cm, minimum height=0.9cm},
  io/.style        = {trapezium, trapezium left angle=70, trapezium right angle=110,
                      draw, minimum width=2.5cm, minimum height=0.9cm},
  processo/.style  = {rectangle, draw, minimum width=2.5cm, minimum height=0.9cm},
  decisione/.style = {diamond, draw, aspect=2, minimum width=2.5cm},
  freccia/.style   = {thick, -{Stealth}}
}

\title{Numeri primi e anello delle classi di resto}
\author{Marco Cerbella}
\date{Lezione 13 -- 2 novembre 2026}

\begin{document}
\maketitle

\section{Conseguenze del teorema di fattorizzazione}

\begin{proposizione}
I numeri primi sono infiniti.
\end{proposizione}

\begin{proof}
Per assurdo siano finiti: $p_1, \dots, p_k$. Consideriamo $m = p_1 p_2 \cdots p_k + 1$. Per ogni $p_i$ si ha $p_i \nmid m$ (il resto della divisione di $m$ per $p_i$ è $1$), ma $m$ ha un divisore primo per il teorema di fattorizzazione: assurdo.
\end{proof}

\begin{proposizione}
Sia $p$ un numero primo ($p > 0$). Allora $\sqrt{p}$ è irrazionale.
\end{proposizione}

\begin{proof}
Dobbiamo dimostrare che $\sqrt{p} \notin \mathbb{Q}$. Per assurdo sia $\sqrt{p} \in \mathbb{Q}$, cioè $\exists a, b \in \mathbb{Z}$, $b \neq 0$, con $\sqrt{p} = \frac{a}{b}$. Allora $b\sqrt{p} = a$ e, elevando al quadrato, $b^2 p = a^2$. Nella fattorizzazione di $a^2$ il primo $p$ ha esponente \emph{pari}; ma nella fattorizzazione di $b^2 p$ il primo $p$ ha esponente \emph{dispari}: assurdo. Quindi $\sqrt{p} \in \mathbb{R}$ ma $\sqrt{p} \notin \mathbb{Q}$.
\end{proof}

\section{$\mathbb{Z}_n$: l'anello delle classi di resto}

Siano $a, b \in \mathbb{Z}$, $n > 0$. Si ha $a \equiv b \bmod n \iff n \mid a - b \iff \exists k \in \mathbb{Z}$ tale che $a - b = nk$.

\[
  \mathbb{Z}_n = \{[0]_n, [1]_n, \dots, [n-1]_n\} = \mathbb{Z}/{\equiv_n}
\]
con
\begin{align*}
  [0]_n &= \{h \in \mathbb{Z} \mid 0 \equiv h \bmod n\} = \{kn \mid k \in \mathbb{Z}\},\\
  [1]_n &= \{h \in \mathbb{Z} \mid 1 \equiv h \bmod n\} = \{kn + 1 \mid k \in \mathbb{Z}\},\\
  &\ \,\vdots\\
  [n-1]_n &= \{h \in \mathbb{Z} \mid n-1 \equiv h \bmod n\} = \{kn + n - 1 \mid k \in \mathbb{Z}\}.
\end{align*}

Definiamo le operazioni $+ : \mathbb{Z}_n\times\mathbb{Z}_n \to \mathbb{Z}_n$ e $\cdot : \mathbb{Z}_n\times\mathbb{Z}_n \to \mathbb{Z}_n$ con
\[
  [a]_n + [b]_n = [a + b]_n, \qquad [a]_n\cdot[b]_n = [a\cdot b]_n .
\]

\begin{proposizione}[Lemma]
Siano $a, b, c, d \in \mathbb{Z}$. Se $a \equiv b \bmod n$ e $c \equiv d \bmod n$ allora
\[
  a + c \equiv b + d \bmod n \qquad\text{e}\qquad a\cdot c \equiv b\cdot d \bmod n .
\]
\end{proposizione}

(Il lemma garantisce che le operazioni sono ben definite, cioè non dipendono dai rappresentanti scelti.)

\begin{proposizione}
$(\mathbb{Z}_n, +, \cdot)$ è un anello commutativo con unità. In particolare $[0]_n$ è l'elemento neutro per $+$ in $\mathbb{Z}_n$; l'opposto di $[a]_n$ è la classe $[-a]_n = [n - a]_n$; $[1]_n$ è l'unità per $\cdot$.
\end{proposizione}

\begin{esempio}[in $\mathbb{Z}_{15}$]
\begin{align*}
  [4]_{15} + [10]_{15} &= [4 + 10]_{15} = [14]_{15},\\
  [10]_{15} + [13]_{15} &= [23]_{15} = [8]_{15} \qquad (23 = 1\cdot 15 + 8),\\
  [3]_{15}\cdot[8]_{15} &= [24]_{15} = [9]_{15} \qquad (24 = 1\cdot 15 + 9).
\end{align*}
\end{esempio}

\subsection{Tavole di $\mathbb{Z}_2$, $\mathbb{Z}_3$, $\mathbb{Z}_4$}

$\mathbb{Z}_2 = \{[0], [1]\}$ (mod 2):
\[
\begin{array}{c|cc} + & 0 & 1 \\ \hline 0 & 0 & 1 \\ 1 & 1 & 0 \end{array}
\qquad
\begin{array}{c|cc} \cdot & 0 & 1 \\ \hline 0 & 0 & 0 \\ 1 & 0 & 1 \end{array}
\]

$\mathbb{Z}_3 = \{[0]_3, [1]_3, [2]_3\}$ (mod 3):
\[
\begin{array}{c|ccc} + & 0 & 1 & 2 \\ \hline 0 & 0 & 1 & 2 \\ 1 & 1 & 2 & 0 \\ 2 & 2 & 0 & 1 \end{array}
\qquad
\begin{array}{c|ccc} \cdot & 0 & 1 & 2 \\ \hline 0 & 0 & 0 & 0 \\ 1 & 0 & 1 & 2 \\ 2 & 0 & 2 & 1 \end{array}
\]
Si ha ad esempio $[2]_3\cdot[2]_3 = [2\cdot 2]_3 = [4]_3 = [1]_3$ (perché $4 = 3\cdot 1 + 1$).

$\mathbb{Z}_4 = \{[0], [1], [2], [3]\}$:
\[
\begin{array}{c|cccc} + & 0 & 1 & 2 & 3 \\ \hline 0 & 0 & 1 & 2 & 3 \\ 1 & 1 & 2 & 3 & 0 \\ 2 & 2 & 3 & 0 & 1 \\ 3 & 3 & 0 & 1 & 2 \end{array}
\qquad
\begin{array}{c|cccc} \cdot & 0 & 1 & 2 & 3 \\ \hline 0 & 0 & 0 & 0 & 0 \\ 1 & 0 & 1 & 2 & 3 \\ 2 & 0 & 2 & 0 & 2 \\ 3 & 0 & 3 & 2 & 1 \end{array}
\]
$\mathbb{Z}_4$ ha divisori dello zero: $[2]_4\cdot[2]_4 = [0]_4$. Inoltre $[3]$ è invertibile in $\mathbb{Z}_4$: $[3]_4\cdot[3]_4 = [1]_4$.

\end{document}
